\documentclass[12pt, a4paper]{article}
\usepackage[utf8]{inputenc}
\usepackage[T2A]{fontenc}
\usepackage[russian]{babel}
\usepackage{geometry}
\geometry{left=2.5cm, right=2.5cm, top=3cm, bottom=2cm}

\usepackage{hyperref}
\usepackage{tocloft}
\usepackage{csquotes}
\usepackage{graphicx}

% Настройка гиперссылок
\hypersetup{
    colorlinks=true,
    linkcolor=blue,
    citecolor=red,
    urlcolor=cyan,
    pdftitle={Название работы},
    pdfauthor={Ваше имя}
}

% ==================== КОМАНДЫ ========================

% Дисклеймер в рамке
\newcommand{\disclaimer}[1]{
    \vspace{1em}
    \noindent\fbox{
        \parbox{\dimexpr\textwidth-2\fboxsep-2\fboxrule}{
            \small\textbf{Дисклеймер:} #1
        }
    }
    \vspace{1em}
}

% Цитаты
\newcommand{\myquote}[1]{
    \textit{\begin{quote}\enquote{#1}\end{quote}}
}

\newcommand{\quoteauthor}[2]{
    \begin{quote}
        \enquote{#1}
        \hfill--- #2
    \end{quote}
}

% ====================== ДОКУМЕНТ ======================

\begin{document}

% Титульный лист
\title{Лекция о вузах}
\author{Константинов Илья Игоревич \\
\href{t.me/iikonstantinov}{telegram: iikonstantinov}}
\date{\today}
\maketitle

\newpage

% ==================== ОГЛАВЛЕНИЕ НА ОТДЕЛЬНОЙ СТРАНИЦЕ ====================
\thispagestyle{empty}           % убираем номер страницы
\begin{center}
    {\Large\bfseries Оглавление}
\end{center}
\vspace{1.5em}

\tableofcontents
\newpage

\newpage

\thispagestyle{empty}
\begin{center}
    {\Large\bfseries Благодарность}
\end{center}
\vspace{1.5em}

Я выражаю глубокую благодарность моим преподавателям Дмитрию Трушину и Марии Горденко, а также моим друзьям и преподавательскому составу ПМИ 2024--2028.

Быть может, вы не помогали мне явным образом, но сказанные вами слова и переданный вами опыт позволили сделать выводы, которые вылились в данный конспект.

\vspace{2em}
\hfill Автор

% ====================== ОСНОВНОЙ ТЕКСТ ======================

\section{Введение}

\subsection{Пролог}

\disclaimer{Данный документ носит исключительно информационный характер и не несёт в себе цели кого-либо оскорбить. Если вы не согласны с высказанным фактом или мнением, контакты автора указаны на титульном листе.

Также текст не имеет никакого отношения ни к одной образовательной организации, с которой связан автор. Материал является исключительно автономной работой, созданной для просвещения.}

Конспект будет руководствоваться принципом, который зачастую принят в математическом сообществе:

\myquote{Какая разница, как ты выглядишь, какими словами говоришь, какие у тебя регалии и достижения, важно лишь то, говоришь ты правду или нет}

\subsection{Базовые понятия}

Когда вам кто-то начинает говорить про жизнь, всегда стоит иметь в виду, кто именно этот человек. Есть большая разница между фактом и мнением: мнение зависит от человека и годно лишь в конкретных обстоятельствах, факт верен вне зависимости от того, кто его произносит. Автор будет пытаться что-то рассказать о жизни, так что периодически будет явно сказано, что какое-либо высказывание является \textit{мнением} автора. В ином случае подразумевается, что написанное является фактом с какой-либо доказательной базой, хотя часто описанные вещи будут иметь исключительно эмпирические \footnote{Страшное слово <<эмпирические>> всюду ниже можно спокойно заменять на <<по жизненному опыту...>>} доказательства, основанные на жизненном опыте, но по каким-то причинам не подвергающиеся сомнению. Иначе было бы оговорено, что это является именно мнением.

Сам конспект посвящён в первую очередь обучению, просто, в связи с обстоятельствами, обучению именно в вузах. В следующем абзаце будут использованы крайне грубые формулировки.

Обучение навыку делится на две части: практическую и теоретическую. Навыком в данном случае может быть что угодно --- как умение писать код, так и умение учиться, жить, что-то доказывать или опровергать. 
\begin{itemize}
    \item Можно ли научиться чему-то, не имея практики? Безусловно, но, когда будете пытаться превратить в реальность заученную теорию, придётся практиковаться и подбивать свои действия под представления о том, <<как это должно быть>>. 
    \item Можно ли научиться чему-то без теории? Да, но тогда всю теоретическую часть будете выводить интуитивно, пытаясь определить правила, которые либо работают, либо не работают. В конечном итоге в любом случае потратите время и на практику, и на теорию, просто где-то больше на одно или другое. 
\end{itemize}

\section{Фундаментальное образование}

\subsection{С чего всё начиналось}

Изначально вуз, НИИ или институт --- это не место образования. В первую очередь это место единения учёных, лабораторий и исследовательской среды ради концентрации научной деятельности. Просто если человек приходит в НИУ ВШЭ как учёный, он обязан взять преподавательские часы, а если приходит как преподаватель, то причислиться к какой-либо лаборатории или, как минимум, базовой кафедре. В первом случае человек науки будет обязан становиться преподавателем и брать на себя эту головную боль, во втором --- никто не заставит тебя писать статьи. В частности, к этому имеет отношение такой вывод: одно из ключевых различий \textit{учителя} и \textit{преподавателя} в том, что учитель зачастую обязан иметь педагогическое образование. Можно очень глубоко и бессмысленно начать копаться в этимологии слов, но важно понять то, что задача преподавателя именно в том, чтобы дать вам возможность учиться, а не насильно чему-то научить. А пользоваться этой возможностью или нет --- уже не их решение.

Но зачем тогда место науки насильно превращается в образовательное учреждение? Просто где-то этих учёных нужно учить. А кто с этим справится лучше, как не другие учёные? Издержкой такого метода является то, что никто не обещает, что хороший учёный будет хорошим преподавателем и сможет вас чему-то научить. Даже на уровне составления образовательных программ часто сидят люди, которые к образованию имеют очень косвенное отношение. И в итоге можно получить программу, созданную идиотами, с курсами, собранными из говна и палок. Поэтому относиться к вузу и администрации как к оракулам, всегда знающим, что и как лучше делать, --- затея очень сомнительная. У вас спокойно может быть предмет, который съест кучу нервов, времени и научит вас чему-то абсолютно бесполезному, причём не просто не имеющему прикладных выводов, а не имеющему пользы даже в любом из вариантов дальнейшего будущего.

Но тем не менее вуз всё-таки стал местом для образования и принял эту участь. К тому же существует немало специальностей, которые требуют огромной горы знаний, которые нигде, кроме как в вузе, не получить. Поэтому перед таким учреждением стояли две задачи --- дать то, что может быть полезно как учёным, так и специалистам. Но ведь и наука, и индустрия постоянно меняются. В любом случае лет через десять ты будешь учить что-то новое. Что тогда давать, чтобы не обновлять программу каждый год? 

\subsection{А что получилось}

В итоге получилось, что ультимативно нужно давать базу, так называемую основу, фундамент, который позволит вам в дальнейшем понять, как работает то, что придумали через двадцать лет после завершения вами института. Один из важнейших навыков --- умение учиться, но про него чуть позже. Получается, задача вуза свелась к обучению основам какой-либо области, что формулируется как \textit{фундаментальное образование}. Это, в свою очередь, вообще не подразумевает необходимость прохождения вуза для того, кто хочет стать грамотным специалистом. Даже тем же программистом можно спокойно стать, не прострадав несколько лет математикой; главное --- знать, какие именно вещи стоит учить. Поэтому одно из важных разделений внутри вуза --- это, грубо говоря, теоретики и практики. Вы будете встречать людей, которые во всю глотку орут, что фундаментальные неприкладные знания --- лишь трата времени, а также других, что будут с горящими глазами сидеть на лекции по биологии, хотя вообще-то планировали стать математиками: <<но ведь полезно же>>. Автор статьи скорее теоретик. Иначе он не смог бы составить такой душный бесполезный текст. Так вот. Теперь мы встретим первую сложность --- точку зрения. Очень много озвученных вещей будут зависеть от вашей точки зрения. Сам конспект предполагает независимость от конкретной позиции, но всегда стоит держать в уме, кто из этих громогласных вам ближе.

\subsection{Что от образования можно взять}

Необходимо идти в университет с чёткой целью, что вы хотите от него в конечном счёте получить: какие именно навыки, какие особенности специальности и так далее. Вы можете спокойно принять, что какой-то предмет вам нафиг не нужен, и поставить его вниз списка приоритетов. А также вы должны понимать, что, скорее всего, ваше мнение изменится до неузнаваемости после первой же сессии. Ключевое, что меняется при поступлении, --- общество. Вокруг вас будут абсолютно незнакомые люди, с разными целями, задачами и опытом. Окружив вас, они покажут и расскажут то, что вам до этого не было возможности узнать. Поэтому цели необходимы, но стоит держать в уме, что эти цели могут поменяться. Жизнь всегда меняется: появляются новые возможности, уходят старые. Но есть какие-то неизменные вещи.

\subsection{А что оно может дать}

Общество --- это одно из главных мест для обучения в университете. Бытует мнение, что школа нам нужна скорее не для знаний, а для социализации и выстраивания отношений в обществе. Это никуда не девается и в дальнейшей жизни. В частности, у этого самого общества вы многому научитесь, может быть, и неявно.

\section{Со стороны студента}

И вот наконец мы подошли к содержательной части лекции. Эгидой дальнейшей части будет фраза: <<Всё гениальное просто>>.

\subsection{Чему учиться?}

Да чему хотите, тому и учитесь. Учите то, что хотите, что не хотите, то не учите. Хотите учить математику --- учите математику. Хотите учить программирование --- учите программирование. В целом можно и заканчивать. 

На самом деле тут можно разделить на две части.

\begin{enumerate}
    \item Предметы, которые учат практическим вещам. 
    
    Тут ключевой материал скорее содержится в заданиях. Зачастую имеет место интересный момент: если вы не слушали лекцию и в заданиях требуется какой-то метод или подход, крайне интересно бывает вывести его самостоятельно. Просто, не зная ничего, попробуйте придумать всё. Чем больше времени вы потратите, тем лучше поймёте, почему представление, которое считается эталонным, выглядит именно так. 
    
    Представим, у нас есть что-то необъятное в виде архитектуры ПО. Если вы, не зная ничего, попробуете написать проект, то потратите невероятное количество времени, переписывая его раз двадцать. Но в этом и смысл. Всегда можно спросить у Гуга, а какое есть готовое решение, но фундаментальность образования заключается именно в том, что это готовое решение на самом деле вы способны придумать самостоятельно.
    \item Предметы, которые не нужны здесь и сейчас. 
    
    Часто студентам технических вузов тяжело даётся вышмат. И очень часто возникает вопрос: а зачем? Вряд ли, когда вы придёте писать микросервис, вас попросят взять интеграл, а ведь деньги дадут только за последнее (почему это неправда, опять же, обсудим позже). 
    
    На самом деле суть математики в этом плане заключается в первую очередь в построении мозгов. Она позволит расшатать ваши мозги так, что после неё многое покажется понятным. Особенно другая математика. Я лично агитирую учить математику и учить её как можно раньше.
\end{enumerate}

Также упомяну про приоритеты. Это правда, что далеко не все предметы \textit{полезные}. Поэтому важно научиться расстановке приоритетов. Оценка --- это, конечно, очень круто, но лучше вы потратите время на профильный предмет, который вам реально пригодится, нежели будете тратить конечный ресурс, время, на что-то вторичное. Так что, когда вы делаете какой-то предмет, нужно понимать, а зачем он реально вам нужен. И может оказаться так, что что-то, что кажется необходимым и тратит много времени, необходимым не является.

\subsection{А как учиться?}

Да как хотите, так и учитесь. В целом можно и заканчивать.

Навык <<учиться>> абсолютно неосязаем. Как понять, что ты умеешь учиться? 

Что выделил тут:

\begin{enumerate}
    \item Конспекты
    
    До меня долго доходило, почему нужны конспекты. На самом деле они нужны не как ваш материал подготовки или дополнительный источник информации. Когда вы приходите на лекцию и ведёте конспекты, работают одновременно три способа запоминания: зрительный, когда смотрите на доску; письменный, когда пишете; и звуковой, когда лектор что-то бубнит. Правда, не всегда получается делать всё вместе качественно, это тоже стоит учесть.
    
    \item Обстановка
    
\begin{quote}
\itshape
\enquote{— Ну как? Обстановка создана? Сейчас начнёшь проверять. \\ — Почему всегда надо издеваться?}

\raggedleft
\textnormal{«Еще раз про любовь», 1968 г.}
\end{quote}

    Многое в эффективной учёбе зависит от обстановки. Если вы не можете сконцентрироваться дома --- учитесь вне дома: существуют коворкинги, библиотеки, кафе. Если не можете учиться одни --- просто скажите об этом в обществе, где люди много учатся. Скорее всего, найдёте компанию, в которой будете проводить время за учёбой вместе. Пункты, которые вы можете изменить:
    \begin{itemize}
        \item Есть люди вокруг / нет людей
        \item Освещение
        \item Музыка
        \item Время
    \end{itemize}
    Я, например, очень продуктивно учусь, когда сижу один в три часа утра в тёмном освещении. Это не стиль. Это продуктивность. 
\end{enumerate}

\subsection{А зачем учиться?}

Тут всё слишком индивидуально. Университет может стать для вас альма-матер, и вы потратите годы жизни на учение. Он может стать для вас просто обязательной галочкой, которую надо закрыть. В нём можно не искать знаний, а лишь общество и времяпрепровождение. Зачастую навык <<учиться>>, а также полученные опыт и общество могут быть полезнее, чем те знания, которые даются на лекции. 

\subsection{Работа}

Огромный холивар создаётся вокруг работы во время учёбы в вузе. Автор на стороне <<не идите работать, пока учитесь>>. Попробую аргументировать свою точку зрения:
\begin{itemize}
    \item Обучение --- это трамплин, работа --- прыжок с этого трамплина. Чем дольше вы учитесь, тем дальше и быстрее прыгнете.
    \item Вуз --- это последнее место, где вы правда будете учиться. Через четыре года вуза вы, скорее всего, окончательно перестанете хотеть учиться. Соответственно, пусть последнее обучение пройдёт максимально качественно.
    \item Время, потраченное на работе вместо учёбы, --- это потерянные деньги в далёком будущем. Зачастую работа гораздо легче учёбы, особенно на высоких должностях. Более того, у каждого студента есть кризис, когда он думает, что уже всё выучил: ведь работа, за которую платят, не так много знаний требует, к тому же они уже есть. Но вот только реальный профит от учёбы вы получите на более высоких должностях, особенно если планируете строить карьеру на hard skills. Тут научное обоснование может взять верх над банальным опытом. 
    \item Единственная причина пойти работать --- вас отчислили из дома. Или вы ушли по собственному желанию. Если у вас есть возможность сидеть на шее родителей в комфорте, то сидите, пожалуйста. Этап сепарации вас от родителей и родителей от вас произойдёт рано или поздно. Абсолютно нормально не работать. Более того, очень многие, например я, мечтают вообще никогда не работать...
    
    Если работа --- это единственное, что даст вам возможность дальше учиться, тогда да, но никогда ради <<опыта>>, <<денег>> и подобного. Выходя из хорошего вуза, вы за несколько лет будете получать столько денег, сколько вообще не думали. Учиться вы будете года четыре, а работать ещё сорок лет минимум.
\end{itemize}

Что делать, если всё-таки по какой-то причине работать надо? 

Нужно, чтобы работа занимала очень мало вашего времени и при этом была достаточно прибыльна. Очень советую пойти преподавать. Преподавателей в хороших школах всегда мало: СУНЦ МГУ, ЛВШЭ, 57, 179, ШЦПМ. Если вы поступили в хороший вуз, значит, вы уже что-то знаете, осталось научиться это преподносить, а это решается опытом. 
Если не хотите преподавать в школах --- занимайтесь репетиторством. Берите 2к/ч, пять занятий в субботу. В итоге, тратя один день в неделю, получаете 40к в месяц. На еду этого хватает, по крайней мере в моей молодости...
Также крайне полезным будет опыт в лабораториях и любой исследовательской деятельности. Никто не заставит вас дальше идти в науку, а в резюме будет смотреться гораздо весомее, чем ничего или преподавание. 

\section{Хороший, плохой, злой}

Любой вуз делится на три столпа:

\begin{itemize}
    \item Студенты
    \item Преподаватели
    \item Администрация
\end{itemize}

Если начать судить вуз, то начинать стоит с этих критериев. Когда говорится <<хороший вуз>>, значит, все три столпа у него работают хорошо. Когда говорится <<плохой вуз>>, значит, хотя бы один из столпов работает плохо. Плохие студенты отпугнут хороших преподавателей. Плохая администрация будет увольнять ни за что и устроит испытание силы воли вместо обучения. Плохой преподаватель может и вовсе отбить желание учиться. Давайте разбирать по порядку.

Я озвучу ряд критериев оценки, не придавая им эмоциональную окраску. Всегда может оказаться, что для одного что-то хорошо, а что-то плохо.

\subsection{Студент}

Если вы говорите: \myquote{Ой, да там препод ничему не научит, так что ботать не буду} --- то вы не правы. Задача учиться --- тоже ваша задача. Хорошо или плохо объясняет преподаватель, не влияет на значимость и красоту предмета. Не можете учиться с этим преподавателем? Ходите к другому. Учитесь по конспектам. Просите друга пересказать лекцию. Ваша задача --- взять всё, что нужно, любыми доступными методами. 

В моих словах можно заметить противоречие. На ком в итоге лежит ответственность за обучение? Правильный ответ: ни на ком. Общение --- это процесс приёма-передачи. И приёмник, и передатчик должны для этого работать хорошо.

Адекватный преподаватель очень любит студентов, задающих вопросы. Да, нужно отличать содержательный вопрос от несодержательного, но, если вы попросите переобъяснить последние пять минут лекции, от этого никто не умрёт. Зато преподаватель поймёт, что именно в его повествовании непонятно, а вы вообще-то пытаетесь его понять. С преподавательской стороны это очень приятно. 

Соответственно, когда мы пытаемся оценить <<студентов>> вуза, то смотрим на итоговые критерии: учатся ли, как строят отношения с остальными, чем в среднем занимаются в свободное время.

\subsection{Преподаватель}

Что такое <<хороший преподаватель>>, всегда зависит от наблюдателя. Но всегда есть простая отсечка: хорошее правительство --- это то, которое хотя бы не мешает. Чуть позже мы поговорим о способах преподавания, когда подойдём к конкретным вузам. Но сейчас выделю такие вещи:
\begin{enumerate}
    \item А умеет ли преподаватель учиться? Меняется ли курс год от года? 
    \item А знает ли преподаватель, кто такие студенты?
    \item А зачем этот человек вообще пошёл преподавать?
\end{enumerate}

Также существует два ортогональных подхода к отношению преподавателя к студенту:
\begin{itemize}
    \item Преподаватель в позиции силы. Преподаватель всегда прав.
    
    Это самая частая позиция в общественной жизни. Она помогает получить тоталитаризм мнения: слово преподавателя есть закон. Значит, ему необходимо верить.
    \item Преподаватель тоже человек, который имеет свойство ошибаться.
    
    Далеко не всем может нравиться такой подход. Но так зачастую можно услышать извинения от преподавателя. Способность признавать ошибки --- это круто.
\end{itemize}

Соответственно, когда мы пытаемся оценить <<преподавателей>> вуза, то смотрим на итоговые критерии: как преподают, зачем преподают, как строят отношения со студентами, какое отношение имеют к предмету.

\subsection{Администрация}

Администрация должна понимать, что вы пришли сюда учиться. В идеале --- как минимум минимально мешать. 

Взаимодействовать с администрацией, например с Учебным офисом, вы будете в лучшем случае. Тут я могу дать такой совет: нет более эффективного метода чего-то добиться от занятого человека, чем прийти лично. Если вам что-то нужно от администрации --- приходите лично и просите сделать.

Также нужно иметь в виду, что есть вещи, которые администрация делать <<не хочет>> и делать <<не может>>. Со вторым вам поможет только бог. С первым всегда можно попробовать чего-то добиться --- например, облегчить жизнь администрации и прийти с уже чем-то готовым. Если на какой-то курс есть отбор, это не значит, что вас не могут туда зачислить, --- просто не хотят, значит, всегда можно что-то придумать. Для примера возьму свой случай взаимодействия с администрацией:
\begin{enumerate}
    \item Начало первого курса. ПМИ. Меня никто не знает, громких достижений нет. Хочу пройти <<Курс по выбору>> второго курса <<Advanced Cpp>>. 
    \item Летом пишу в личку лектору с вопросом, можно ли. Получаю положительный ответ и просьбу подойти в начале года.
    \item Пишу в начале года, получаю игнор.
    \item В начале года подхожу лично после лекции. Получаю подтверждение одобрения.
    \item Пишу академическому руководителю с отсылкой на одобрение лектора. Узнаю, что его вообще можно взять вместо обычного курса <<плюсов>>. Получаю положительный ответ.
    \item Прихожу к администрации со всеми картами на руках: одобрение от лектора, одобрение от академического руководителя. Задачей администрации остаётся только перевести меня на курс. 
\end{enumerate}
Таким образом, на первой же неделе учёбы я уже изменил себе академический план, хотя этих людей (и они меня) видел впервые в жизни. 

Также ортогональный случай. 
\begin{enumerate}
    \item Начало второго курса. Хочу пойти на доступный курс по выбору <<Распределённые системы>>, который для многих является обязательным на третьем курсе.
    \item Подаю заявку в выборе КПВ на это. 
    \item Через месяц смен КПВ мой статус остаётся непонятным, никто не даёт знать, перешёл я или нет.
    \item Пишу лично преподавателю. Он говорит, что мне не хватает пререквизитов.
    \item В последний вагон пишу заявление на перевод на другой курс, куда меня зачисляют.
    \item В итоге --- потерянный месяц обучения. А можно было просто спросить...
\end{enumerate}

\section{Со стороны преподавателя}

Тут я вставил блоки, которые связаны с преподаванием. Для студентов они скорее помогут понять какие-то механики происходящего. И точно определят, когда что-то идёт не так.

\subsection{Вред презентаций}

Презентации в образовании --- это зло, причём крайне редко оправданное. В преподавателе и студенте относительно предмета ключевая разница --- скорость восприятия информации. Когда преподаватель пишет на доске, скорость его письма совпадает со скоростью восприятия студента. Когда листаются слайды, эта связь пропадает. В итоге один перелистнутый слайд --- и преподаватель уже не понимает, на каком моменте осознания его студенты. Ведь преподаватель знает заранее содержание слайда, а студент тратит время на его осознание. Тогда студент будет сначала пять минут читать слайд, пока преподаватель говорит какую-то содержательную информацию. Если для восприятия этой информации не нужен слайд, тогда он просто тратит силы и фокус студента. Если содержательную информацию от преподавателя можно пропустить, тогда зачем нужен преподаватель?

Никто за последние четыре миллиарда лет не придумал ничего лучше, чем маркер и доска. Ну или мел. 

\subsection{Распределение по уровням}

Сейчас будет немного душноты. 

Положим, вы уже знаете, что такое \href{https://ru.wikipedia.org/wiki/%D0%9D%D0%BE%D1%80%D0%BC%D0%B0%D0%BB%D1%8C%D0%BD%D0%BE%D0%B5_%D1%80%D0%B0%D1%81%D0%BF%D1%80%D0%B5%D0%B4%D0%B5%D0%BB%D0%B5%D0%BD%D0%B8%D0%B5}{нормальное распределение}.

\centering
\includegraphics[scale=0.4]{./.img/graph.png}

Если мы попытаемся взять группу студентов в 30 человек и измерить их по какому-либо понятию <<умности>>, то оно и получится. Ну либо <<логарифмическое нормальное>>, но не суть. 

Преподаватель может охватить какой-то диапазон студентов по этому понятию <<умности>>, причём этот параметр измеряется в абсолютных измерениях, то есть он работает с отрезком студентов $[x, y]$. Но вот только нормальное распределение зависит от двух параметров --- среднего и дисперсии. И когда дисперсия слишком большая или среднее не там, большой процент студентов оказывается <<за бортом>>. Как это называется? Несбалансированная группа. В своё время я такое встретил в ЛВШЭ, когда преподавал там программирование. Например, 40\% студентов, которые могли взять гораздо больший материал, не могли его получить, потому что я должен был охватить остальные 60. Именно тут проявляется проблема большого потока. Чем меньше группа, тем более качественное понимание материала получает каждый, потому что преподаватель лучше понимает актуальные проблемы и задачи. Оптимальная семинарская группа --- 15 человек, жаль, не каждый может себе такое позволить. 

Мысль эта к тому, что у хорошего профильного предмета должно быть деление по уровню. Это не затем, чтобы выделить элитарных студентов, --- совсем нет. Если преподаватель понимает, что у него <<слабая группа>>, то он подстроит материал под неё, и вследствие этого повысится обучаемость. А <<умненьких>>, возможно, вообще стоит начать переучивать, потому что там кто во что горазд.

\end{document}